\begin{landscape}
\begingroup
\fontsize{10}{13}\selectfont
\fancyhead[L]{\small NACHLASS. STÖRUNGEN DER PALLAS}
\renewcommand{\headrulewidth}{0pt}
\par\medskip
\SourceLeaf{563}
\pdfbookmark[2]{Nachlass VI [9]}{pallas-VI-fragment9}
\begin{center}[9.]\end{center}
\makeatletter
\providecommand{\VInineHistoricalNote}[1]{\begingroup
\renewcommand{\thefootnote}{*)}
\long\def\@makefntext##1{\noindent [*)\ ##1}
\footnote{#1}\endgroup}
\makeatother
\par\noindent\hspace*{\parindent}[In Artikel 19 der Exposition wurde bereits bemerkt, dass Pallas einen Fall bietet, in dem die mittleren Bewegungen \(n\) und \(n'\) in einem rationalen Verhältnis stehen. Die mittleren Bewegungen von Pallas und Jupiter verhalten sich nämlich genau wie \(18:7\), wobei sich dieses Verhältnis stets genau wiederherstellt\VInineHistoricalNote{Vgl. \textsc{Gauss} an \textsc{Bessel}, 5. Mai 1812, S. 421 dieses Bandes und Göttingische gelehrte Anzeigen vom 25. April 1812, Bd. VI, S. 349.]}. In diesem Artikel sollen die vom Argument \(18\text{\mdseries\AstroSymbols ♃}-7\text{\mdseries\AstroSymbols ⚴}\) abhängigen Glieder für die praktische Rechnung in säkularer Form angegeben werden. Der nächste Artikel wird die exakte Integration behandeln, durch die die periodische Form wiederhergestellt wird.
\par\noindent\hspace*{\parindent}Aus den obigen numerischen Werten für \(T,V,W\) erhält man nach der Interpolation wie in den Artikeln 2—7 die vom genannten Argument abhängigen Glieder
\[
\begin{aligned}
\text{in }\frac{(1-\cos i)\mathrm{d}\AscendingNode}{n\mathrm{d}t}:&\quad
0\overset{\prime\prime}{,}000\cos(18M'-7M)\\
&\qquad+0\overset{\prime\prime}{,}001\sin(18M'-7M)\\[.7ex]
\text{in }\frac{\mathrm{d}\tilde\omega}{n\mathrm{d}t}-\frac{(1-\cos i)\mathrm{d}\AscendingNode}{n\mathrm{d}t}:&\quad
+0\overset{\prime\prime}{,}095\cos(18M'-7M)\\
&\qquad-0\overset{\prime\prime}{,}055\sin(18M'-7M)\\[.7ex]
\text{in }\frac{\mathrm{d}\log\mathrm{hyp}\,n}{n\mathrm{d}t}=-\tfrac32\frac{\mathrm{d}\log\mathrm{hyp}\,a}{n\mathrm{d}t}:&\quad
+0\overset{\prime\prime}{,}012\cos(18M'-7M)\\
&\qquad-0\overset{\prime\prime}{,}027\sin(18M'-7M)\\[.7ex]
\text{in }\frac{\mathrm{d}\varepsilon}{n\mathrm{d}t}-\frac{(1-\cos i)\mathrm{d}\AscendingNode}{n\mathrm{d}t}:&\quad
-0\overset{\prime\prime}{,}039\cos(18M'-7M)\\
&\qquad-0\overset{\prime\prime}{,}021\sin(18M'-7M).
\end{aligned}
\]
\par\noindent\hspace*{\parindent}Fasst man die beiden Glieder in \(\cos(18M'-7M)\) und \(\sin(18M'-7M)\) zu einem zusammen und führt die mittleren Längen über die Werte
\[
\begin{aligned}
M'&=\text{\mdseries\AstroSymbols ♃}-11^0 17'\quad5\overset{\prime\prime}{,}39\\
M&=\text{\mdseries\AstroSymbols ⚴}-121\quad8\quad54{,}50,
\end{aligned}
\]
ein, so erhält man das Glied
\[
\begin{aligned}
\text{in }\tilde\omega:&\quad-0\overset{\prime\prime}{,}1093\int\cos(18\text{\mdseries\AstroSymbols ♃}-7\text{\mdseries\AstroSymbols ⚴}+134^0 31'\,38'')\mathrm{d}\text{\mdseries\AstroSymbols ⚴}\\
\text{in }\log a:&\quad+0{,}415\int\cos(18\text{\mdseries\AstroSymbols ♃}-7\text{\mdseries\AstroSymbols ⚴}+170^0 56'\,57'')\mathrm{d}\text{\mdseries\AstroSymbols ⚴}\\
&\qquad\text{(in Einheiten der 7. Dezimalstelle)}\\
\text{in }\int n\mathrm{d}t:&\quad-0\overset{\prime\prime}{,}02955\iint\cos(18\text{\mdseries\AstroSymbols ♃}-7\text{\mdseries\AstroSymbols ⚴}+170^0 56'\,57'')\mathrm{d}\text{\mdseries\AstroSymbols ⚴}^{2}\\
\text{in }\varepsilon:&\quad-0\overset{\prime\prime}{,}0438\int\cos(18\text{\mdseries\AstroSymbols ♃}-7\text{\mdseries\AstroSymbols ⚴}+257^0 45'\,45'')\mathrm{d}\text{\mdseries\AstroSymbols ⚴};
\end{aligned}
\]
Die Integration ergibt, wenn man für \(18\text{\mdseries\AstroSymbols ♃}-7\text{\mdseries\AstroSymbols ⚴}\) den konstanten Wert \(125^0 36'\,49''\) einsetzt:
\begin{center}
\begin{tabular}{@{}cccll@{}}
Säkularänderung&von&\(\tilde\omega\)&\(:+0\overset{\prime\prime}{,}0255\)&(jährlich)\\
»&»&\(\log a\)&\(:+0{,}25\)&»\\
&&\multicolumn{3}{l}{(in Einheiten der 7. Dezimalstelle)}\\
»&»&\(\varepsilon\)&\(:-0\overset{\prime\prime}{,}0548\)&»\\
\multicolumn{3}{r}{Faktor von \(tt\) in \(\int n\mathrm{d}t\)}&\(:-0\overset{\prime\prime}{,}0123\)&»
\end{tabular}
\end{center}
\par\noindent\hspace*{\parindent}Für die Neigung und die Länge des Knotens sind die entsprechenden Glieder sehr klein.]
\clearpage
\endgroup
\end{landscape}
